\section*{الفصل \ref{ch:g10:numbers} --- الأعداد ومجموعات الأعداد}

\begin{solution}{exo:g10:numbers:1}
$\frac{15}{3} = 5 \in \N$. \quad $-7 \in \Z$. \quad
$\frac{22}{7} \in \Q$ (ليس عددًا صحيحًا: $22 = 7 \times 3 + 1$).
\quad $\sqrt 9 = 3 \in \N$. \quad $\sqrt{10} \in \R$ (أصمّ، لأن
$10$ ليس مربع \omterm{def:g10:numbers:sets}{عدد ناطق} --- ونقبل ذلك هنا، على منوال
\cref{thm:g10:numbers:sqrt2}). \quad $-2.4 = -\frac{24}{10} \in \Q$.
\quad $0 \in \N$.
\end{solution}

\begin{solution}{exo:g10:numbers:2}
(أ) $\intco{-2}{5}$: نقطة مملوءة عند $-2$، ونقطة مجوّفة عند $5$.
(ب) $\intoo{3}{+\infty}$: نقطة مجوّفة عند $3$، وتظليل نحو اليمين.
(ج) $\intoc{-\infty}{-1}$: تظليل من اليسار حتى نقطة مملوءة عند
$-1$.
(د) ``المسافة من $x$ إلى $2$ لا تتجاوز $3$'' تعني $\abs{x - 2} \leq 3$،
أي $x \in \intcc{-1}{5}$: نقطتان مملوءتان عند $-1$ و $5$.
\end{solution}

\begin{solution}{exo:g10:numbers:3}
(أ) تتراكب الفترتان بين $0$ و $2$:
$I \cap J = \intcc{0}{2}$ (لأن $0 \in J$ و $2 \in I$، وكلاهما ينتمي إلى
المجموعة الأخرى أيضًا)، و$I \cup J = \intco{-3}{4}$.

(ب) $I \cap J$ هي مجموعة الأعداد $x$ التي تحقق $-2 \leq x$ و $x < 1$:
$\intco{-2}{1}$. \omterm{def:g10:numbers:interunion}{والاتحاد} يغطي كل شيء: $I \cup J = \R$.
\end{solution}

\begin{solution}{exo:g10:numbers:4}
$\abs{-6} = 6$؛ \quad $\abs{4 - 9} = \abs{-5} = 5$؛ \quad
$\abs{-3 - 5} = \abs{-8} = 8$؛ \quad
$\sqrt 2 > 1$ إذن $\abs{\sqrt2 - 1} = \sqrt2 - 1$؛ \quad
$1 - \sqrt2 < 0$ إذن $\abs{1 - \sqrt2} = \sqrt2 - 1$ أيضًا (فللعدد
ولمقابله \omterm{def:g10:numbers:abs}{القيمة المطلقة} نفسها).
\end{solution}

\begin{solution}{exo:g10:numbers:5}
$\abs{x} = 5$: المسافة $5$ من $0$، إذن $x = 5$ أو $x = -5$.

$\abs{x - 1} = 3$: المسافة $3$ من $1$، إذن $x = 4$ أو $x = -2$.

$\abs{x - 4} \leq 1$: المسافة لا تتجاوز $1$ من $4$، إذن
$x \in \intcc{3}{5}$.

$\abs{x + 2} < 3$: المسافة أصغر تمامًا من $3$ من $-2$، إذن
$x \in \intoo{-5}{1}$.
\end{solution}

\begin{solution}{exo:g10:numbers:6}
توصف كل \omterm{def:g10:numbers:interval}{فترة} بمركزها $a$ (منتصف الطرفين) وبنصف قطرها $r$
(نصف الطول).

$\intcc{2}{8}$: $a = 5$ و $r = 3$، إذن $\abs{x - 5} \leq 3$.

$\intoo{-5}{1}$: $a = -2$ و $r = 3$، والطرفان مفتوحان، إذن $\abs{x + 2} < 3$.

$\intcc{-7}{-3}$: $a = -5$ و $r = 2$، إذن $\abs{x + 5} \leq 2$.
\end{solution}

\begin{solution}{exo:g10:numbers:7}
لنضع $x = 0.272727\dots$ عندئذٍ $100x = 27.2727\dots$، وبالطرح:
\[
100x - x = 27.2727\dots - 0.2727\dots = 27,
\]
إذن $99x = 27$ و $x = \frac{27}{99} = \frac{3}{11}$، وهو خارج قسمة
عددين صحيحين: أي أن $x$ ناطق.
\end{solution}

\begin{solution}{exo:g10:numbers:8}
\emph{1. صحيح:} جمع عددين صحيحين (موجبين أو سالبين)
يعطي دائمًا عددًا صحيحًا.

\emph{2. خطأ:} $\frac{1}{2}$ خارج قسمة العددين الصحيحين $1$ و $2$
وليس عددًا صحيحًا.

\emph{3. صحيح:} $\frac pq + \frac{p'}{q'} = \frac{pq' + p'q}{qq'}$ هو
من جديد خارج قسمة عددين صحيحين (بمقام غير معدوم).

\emph{4. صحيح:} لنفترض أن $r$ ناطق و $t$ أصمّ وأن $r + t = s$
ناطق. عندئذٍ يكون $t = s - r$ فرقًا بين عددين ناطقين،
أي ناطقًا (حسب 3، مطبقة على $-r$) --- وهذا تناقض. إذن $r + t$
أصمّ.
\end{solution}

\begin{solution}{exo:g10:numbers:9}
بضرب $1.414 \leq \sqrt2 \leq 1.415$ في $2 > 0$:
$2.828 \leq 2\sqrt2 \leq 2.830$.

وبإضافة $3$: $4.414 \leq \sqrt2 + 3 \leq 4.415$.

وبالضرب في $-1 < 0$ ينقلب اتجاه المتراجحتين:
$-1.415 \leq -\sqrt2 \leq -1.414$.
\end{solution}

\begin{solution}{exo:g10:numbers:10}
لنبدأ بالواقعة المساعدة. اقسم $p$ على $3$: الباقي $0$ أو $1$ أو
$2$، أي $p = 3k$ أو $p = 3k+1$ أو $p = 3k+2$. وبالتربيع:
\[
(3k)^2 = 3(3k^2), \quad
(3k+1)^2 = 3(3k^2 + 2k) + 1, \quad
(3k+2)^2 = 3(3k^2 + 4k + 1) + 1 .
\]
والأول وحده مضاعف للعدد $3$: فإذا كان $p^2$ مضاعفًا للعدد $3$، فإن
$p$ كذلك.

لنفترض الآن أن $\sqrt3 = \frac pq$ كسر مختزل تمامًا. بالتربيع نجد
$p^2 = 3q^2$، إذن $p^2$ مضاعف للعدد $3$، إذن $p = 3k$. عندئذٍ
$9k^2 = 3q^2$، إذن $q^2 = 3k^2$ ويكون $q$ مضاعفًا للعدد $3$ أيضًا ---
لكن الكسر $\frac pq$ ما كان إذن مختزلًا تمامًا: تناقض. إذن
$\sqrt3$ عدد أصمّ.
\end{solution}

\begin{solution}{pb:g10:numbers:1}
\textbf{1.} $-7 \in \Z$؛ \ $\frac{13}{4} \in \Q$؛ \
$\sqrt{16} = 4 \in \N$؛ \ $0.121212\ldots = \frac{12}{99} =
\frac{4}{33} \in \Q$ (مسألة نهاية الأسبوع عن العشريات الدورية في الكتاب السابق)؛ \
$\sqrt8 = 2\sqrt2$ عدد أصمّ: أصغر مجموعة تحتويه $\R$؛ \ $\pi$:
$\R$.

\textbf{2.} $\frac pq + \frac rs = \frac{ps + qr}{qs}$ و
$\frac pq \times \frac rs = \frac{pr}{qs}$: خارجا قسمة
عددين صحيحين بمقام غير معدوم --- إذن ناطقان.

\textbf{3.} (أ) لو كان $r$ ناطقًا و $x$ أصمّ و
$r + x = q$ ناطقًا، لكان $x = q - r$ فرقًا
بين عددين ناطقين، أي ناطقًا (السؤال 2):
تناقض. (ب) إذا كان $r \neq 0$ و $r x = q$ ناطقًا، فإن
$x = \frac qr$ ناطق: تناقض.

\textbf{4.} العددان $\sqrt2$ و $-\sqrt2$ أصمّان ومجموعهما $0$؛
والعددان $\sqrt2$ و $\sqrt2$ جداؤهما $2$. نتائج ناطقة من
مكوّنات \omterm{ex:g10:numbers:classify}{صمّاء}: فالأعداد \omterm{ex:g10:numbers:classify}{الصمّاء} لا تشكّل مملكة مستقرة.

\textbf{5.} $\sqrt8 = 2\sqrt2$، إذن
$\sqrt2 + \sqrt8 = 3\sqrt2$: أصمّ (السؤال 3ب، بمعامل
$3$). و$\sqrt2 \times \sqrt8 = \sqrt{16} = 4 \in \N$: جداء
عددين أصمّين يحطّ في أصغر الممالك
جميعًا.

\textbf{6.} $3.475$، ثم $3.471$ (أو $3.4701$ أو $3.47001$
\dots): إلحاق أرقام عشرية ينتج أعدادًا ناطقة جديدة
بين العددين بلا نهاية.

\textbf{7.} تسير مضاعفات $10^{-n}$ على المستقيم بخطوة
مقدارها $10^{-n} < g$. وأول مضاعف أكبر تمامًا
من $a$ --- وهو موجود، لأن المضاعفات تتجاوز $a$ في
النهاية --- يقع على بعد خطوة واحدة على الأكثر من $a$، أي قبل
$a + g = b$: فهو محصور تمامًا بين $a$ و $b$. ونقطة الشبكة
عدد عشري، أي \omterm{def:g10:numbers:sets}{عدد ناطق}: فكل \omterm{def:g10:numbers:interval}{فترة} ذات طول موجب
تحوي عددًا ناطقًا.

\textbf{8.} بين $a$ \omterm{def:g10:numbers:sets}{والعدد الناطق} $r_1$ الموجود في
السؤال 7 يوجد (بالسؤال 7 مرة أخرى) \omterm{def:g10:numbers:sets}{عدد ناطق} $r_2$؛ وبين
$a$ و $r_2$ \omterm{def:g10:numbers:sets}{عدد ناطق} $r_3$؛ وهكذا: عدد لا نهائي منها، كلها
مختلفة، وكلها في \omterm{def:g10:numbers:interval}{الفترة} الأصلية.

\textbf{9.} $r + \frac{\sqrt2}{10^k}$ مجموع \omterm{def:g10:numbers:sets}{عدد ناطق}
وعدد أصمّ (فالعدد $\frac{\sqrt2}{10^k}$ أصمّ حسب
السؤال 3ب)، إذن هو أصمّ. وبما أن
$\frac{\sqrt2}{10^k} < \frac{2}{10^k}$ ينكمش تحت أي حد،
فإنه من أجل $k$ كبير بما يكفي يبقى $r + \frac{\sqrt2}{10^k}$
قبل $b$: أي عدد أصمّ داخل \omterm{def:g10:numbers:interval}{الفترة} --- وبتغيير
$k$ نحصل على عدد لا نهائي منها.

\textbf{10.} لا يوجد أصغر \omterm{def:g10:numbers:sets}{عدد حقيقي} موجب: لو كان $s > 0$ هو
ذاك، لاحتوت \omterm{def:g10:numbers:interval}{الفترة} $(0, s)$ عددًا ناطقًا رغم ذلك (السؤال 7)،
موجبًا وأصغر من $s$. ولا يوجد عدد يأتي مباشرة بعد $3$: فأي
مرشح $c > 3$ يترك \omterm{def:g10:numbers:interval}{الفترة} $(3, c)$، وهي ليست
خالية --- إنها لعبة مسألة نهاية الأسبوع عن انعدام العدد التالي في الكتاب السابق، وقد صارت مبرهنة عن
$\R$.

\textbf{11.} $\abs{x - 5} = 2$: $x = 3$ أو $x = 7$.
$\abs{x - 5} < 2$: $x \in \intoo{3}{7}$.
$\abs{x + 1} \geq 3$: المسافة إلى $-1$ لا تقل عن $3$:
$x \in \intoc{-\infty}{-4} \cup \intco{2}{+\infty}$.

\textbf{12.} $\abs{d - 80} \leq 0.05$، أي
$d \in \intcc{79.95}{80.05}$. المكبس ذو القياس $79.97$ يُقبل؛
أما ذو القياس $80.06$ فيُرفض بفارق جزء من مئة من المليمتر.

\textbf{13.} في حالة $(3, -5)$: $\abs{-2} = 2 \leq 8$. وفي حالة $(-2, -7)$:
$\abs{-9} = 9 = 2 + 7$: مساواة. عند تساوي الإشارتين: $\abs{a + b}$ هو
مجموع المسافتين، ويساوي $\abs a + \abs b$. وعند اختلاف
الإشارتين: يرتدّ المسار على نفسه، فيكون $\abs{a + b}$
\emph{فرق} المسافتين، وهو أصغر تمامًا من
مجموعهما (إلا إذا كانت إحداهما $0$). وتتحقق المساواة بالضبط حين يكون $a$ و
$b$ من الإشارة نفسها أو حين يكون أحدهما معدومًا.

\textbf{14.} الحلول هي النقط المتساوية البعد عن $2$
وعن $-4$: أي المنتصف، $x = -1$. إنه محور
القطعة في مسألة نهاية الأسبوع عن المرآتين في الكتاب السابق، منهارًا إلى بعد
واحد: نقطة وحيدة.

\textbf{15.} $\sqrt{x^2} = \abs{x}$، لا $x$: فمن أجل $x = -3$ يكون
$\sqrt{9} = 3 = \abs{-3} \neq -3$. عندئذٍ تُقرأ $x^2 < 9$
$\sqrt{x^2} < 3$، أي $\abs x < 3$:
$x \in \intoo{-3}{3}$.

\textbf{16.} لا يؤكد القياس سوى أن الطول يقع
في \omterm{def:g10:numbers:interval}{الفترة} $\intcc{1.41420}{1.41422}$ --- وبالسؤالين
7 و 9 \emph{تحوي تلك \omterm{def:g10:numbers:interval}{الفترة} عددًا لا نهائيًا من \omterm{def:g10:numbers:sets}{الأعداد الناطقة}
وعددًا لا نهائيًا من \omterm{ex:g10:numbers:classify}{الصمّاء}}. والأمر نفسه يصح مع كل
سماحية مقبلة، مهما صغرت. فما من قياس يستطيع أن
يفصل بين العائلتين؛ ولا يحسم ذلك إلا برهان عن الطول بالضبط
--- كالبرهان الخاص بقطر مربع الوحدة
(مسألة نهاية الأسبوع عن الصمم في الكتاب السابق).

\textbf{17.} إنها تقترب من $\sqrt2$ بأخطاء أصغر من
$10^{-1}, 10^{-2}, 10^{-3}, \dots$: \omterm{def:g10:numbers:sets}{فالأعداد الناطقة} تلتصق بالعدد
$\sqrt2$ في كل سلّم. والواقعة العامة: كل \omterm{def:g10:numbers:sets}{عدد حقيقي}
تقترب منه \omterm{def:g10:numbers:sets}{الأعداد الناطقة} (بتوره العشرية) بأي قدر
مطلوب --- إنها الكثافة مرة أخرى، منظورًا إليها من جهة
الهدف.

\textbf{18.} $a + b \in \intcc{2.4 + 1.1}{2.6 + 1.3} =
\intcc{3.5}{3.9}$: ارتياب $\pm 0.2$؛ تُجمع الأخطاء.
$a \times b \in \intcc{2.4 \times 1.1}{2.6 \times 1.3} =
\intcc{2.64}{3.38}$: ارتياب مقداره نحو $\pm 0.37$ حول
$3$ --- فالأخطاء النسبية تتجمع في الجداءات، ولذلك تتدهور الدقة
أسرع.

\textbf{19.} اكتب القيمتين الحقيقيتين $a = 2.5 + e_1$ و
$b = 1.2 + e_2$ حيث $\abs{e_1}, \abs{e_2} \leq 0.1$. عندئذٍ يكون
خطأ المجموع $\abs{e_1 + e_2} \leq \abs{e_1} +
\abs{e_2} \leq 0.2$: فقاعدة المهندس \emph{هي} متراجحة
المثلث في السؤال 13.

\textbf{20.} مستقيم الأعداد بلا ثقوب وبلا جيران: فبعد
أي عدد، لا يوجد عدد ``تالٍ'' (السؤال 10). وتنقسم نقطه
إلى أعداد ناطقة --- وهي بالضبط الأعداد العشرية المتوقفة أو
الدورية (مسألة نهاية الأسبوع عن العشريات الدورية في الكتاب السابق) --- وأعداد \omterm{ex:g10:numbers:classify}{صمّاء}،
والعائلتان متشابكتان تشابكًا وثيقًا إلى حدّ أن كل
\omterm{def:g10:numbers:interval}{فترة} تحوي عددًا لا نهائيًا من كل منهما (السؤالان 8 و 9)،
ولذلك لا يستطيع أي قياس، بل البرهان وحده، أن يميّز بينهما
(السؤال 16). وتبقى المفاجأة الأخيرة وعدًا مؤجلًا:
الأعداد \omterm{ex:g10:numbers:classify}{الصمّاء} أكثر عددًا من الناطقة --- لا بالعدّ واحدًا
واحدًا، بل بالمعنى الدقيق لمقارنة اللانهايات، وهي
نظرية تُبنى في الكتب الجامعية.
\end{solution}
